\documentclass[11pt,a4paper]{article}
\usepackage[utf8]{inputenc}
\usepackage[german]{babel}
\usepackage{amsmath,amssymb,amsfonts}
\usepackage{booktabs}
\usepackage{geometry}
\usepackage{xcolor}
\usepackage{hyperref}
\usepackage{graphicx}
\usepackage{listings}

\geometry{margin=1in}

% Projektdaten
\title{\textbf{Report: Generative Frameworks \\ Evaluation und Performance-Metriken}}
\author{Joachim Hofmann}
\date{Juni 2026}

\begin{document}

\maketitle

\section{Introduction}
Dieses Forschungsdokument evaluiert und vergleicht die Leistungseigenschaften von drei modernen generativen Frameworks: \textit{Denoising Diffusion Probabilistic Models} (DDPM), \textit{Flow-Matching} (FM) und \textit{Rectified Flow} (RF). Der Fokus liegt dabei auf der Balance zwischen Berechnungsaufwand und mathematischer Präzision bei der Trajektorienführung.

\section{Einfache Definition der beiden Forschungsziele}
Um die Leistung und Effizienz generativer Frameworks objektiv zu bewerten, orientiert sich diese Untersuchung an zwei zentralen Forschungszielen:

\subsection{Ziel 1: Minimierung der operationalen Kosten (\textit{Wall-Clock Cost} / Inferenz-Effizienz)}
\textbf{In einfachen Worten:} Dieses Ziel befasst sich mit der Geschwindigkeit, mit der ein fertig trainiertes Modell ein neues Bild erzeugen kann. Selbst wenn verschiedene Modelle eine identische neuronale Kernstruktur (z.~B. dasselbe U-Net) nutzen, benötigen manche mathematische Formulierungen deutlich weniger iterative Schleifen (Rechenschritte), um zu einem sauberen Ergebnis zu gelangen.

\textbf{Abkürzungen in der Ergebnistabelle:}
\begin{itemize}
    \item \textbf{NFE} (\textit{Number of Function Evaluations}): Die Gesamtzahl der separaten Aufrufe, die das neuronale Netzwerk durchlaufen muss, um ein einziges finales Bild aus dem Rauschen zu konstruieren.
    \item \textbf{Inferenz-Zeit (ms)} (\textit{Latency}): Die real gemessene physikalische Zeit in Millisekunden, die die Hardware benötigt, um den gesamten Generierungszyklus abzuschließen.
\end{itemize}

\subsection{Ziel 2: Stabilität der mathematischen Formulierung (\textit{Stability} / Geometrie \& Pfad-Präzision)}
\textbf{In einfachen Worten:} Klassische Diffusionsmodelle sind mathematisch hochgradig volatil. Sie sind darauf angewiesen, dass ein künstlicher Rausch-Zeitplan (\textit{Variance Schedule}) exakt kalibriert ist. Ist dieser Zeitplan minimal fehlerhaft, kollabiert das Modell. 

Dieses Forschungsziel soll zeigen, dass der Verzicht auf komplexe Zeitpläne zugunsten einer direkten Pfadbegradigung (\textit{Least Squares}) zu einer extrem stabilen und linearen Lernkurve führt.

\textbf{Abkürzung in der Ergebnistabelle:}
\begin{itemize}
    \item \textbf{Abweichung (MSE)} (\textit{Mean Squared Error}): Der mittlere quadratische Fehler. Er misst, wie stark der vom Modell vorhergesagte Pfad von einer perfekt geraden Linie zwischen dem Startrauschen und dem Zielbild abweicht. Ein stabiler, gerader Pfad minimiert Integrationsfehler im Solver.
\end{itemize}

\section{Die Rolle des MLP in der Zeit-Konditionierung}
Obwohl das strukturelle Rückgrat für die Bildverarbeitung aus einem zweidimensionalen, konvolutionalen U-Net besteht, spielt ein \textbf{Multi-Layer Perceptron (MLP)} eine entscheidende Rolle bei der Einspeisung der zeitlichen Dynamik.

Generative Fluss- und Diffusionsmodelle müssen zu jedem Zeitpunkt wissen, an welcher Stelle der Generierungs-Zeitachse (von purem Rauschen $t=0$ bis zum scharfen Bild $t=1$) sie sich befinden. Dieser Prozess läuft wie folgt ab:
\begin{enumerate}
    \item Die skalare Zeitvariable $t$ wird mittels sinusoidalen Positions-Encodings in einen hochdimensionalen Vektor transformiert.
    \item Ein \textbf{MLP}, bestehend aus linearen Schichten und nicht-linearen Aktivierungsfunktionen ($\text{nn.Sequential}(\text{nn.Linear}(\dots), \text{nn.ReLU}())$), verarbeitet diesen Einbettungsvektor.
    \item Der Ausgang des MLP wird auf die passenden Kanal-Dimensionen projiziert und direkt zu den verborgenen Feature-Maps der Convolutional-Layer addiert. Dadurch wird das Verhalten des U-Nets dynamisch anhand des aktuellen Zeitschritts gesteuert.
\end{enumerate}

\section{Mathematische Algorithmen}

\subsection{Algorithmus 1: Training von 1-Rectified Flow (Rektifizierung)}
Dieser Algorithmus trainiert das neuronale Basis-Netzwerk zur Parametrisierung eines Geschwindigkeitsfeldes $v_\theta$, welches der Richtung linearer Pfade zwischen den Verteilungen $\pi_0$ (Quellrauschen) und $\pi_1$ (Ziel-Datenverteilung) folgt.

\textbf{Ablauf (Bis zur Konvergenz wiederholen):}
\begin{enumerate}
    \item Ziehe ein unabhängiges Paar $(X_0, X_1)$ aus der gemeinsamen Kopplung $\pi_0 \times \pi_1$.
    \item Ziehe einen zufälligen, kontinuierlichen Zeitschritt: $t \sim \text{Uniform}([0, 1])$.
    \item Berechne die lineare Interpolation (den Fluss-Pfad) zum Zeitpunkt $t$:
    \begin{equation}
        X_t = t X_1 + (1 - t) X_0
    \end{equation}
    \item Aktualisiere das Modell $v_\theta$ durch Minimierung des unbeschränkten Least-Squares-Problems via Gradientenabstieg:
    \begin{equation}
        \min_\theta \mathbb{E}_{X_0, X_1, t} \left[ \| (X_1 - X_0) - v_\theta(X_t, t) \|^2 \right]
    \end{equation}
\end{enumerate}
\textbf{Output:} Das gelernte 1-Rectified Flow Vektorfeld $v_\theta$.

\subsection{Algorithmus 2: Inferenz (Sampling / ODE-Integration)}
Um ein neues Bild aus reinem Rauschen zu generieren, wird die gelernte gewöhnliche Differentialgleichung (ODE) vorwärts von $t=0$ bis $t=1$ gelöst.

\textbf{Ablauf via Euler-Integration über $N$ Schritte:}
\begin{enumerate}
    \item Sample Startrauschen aus der Quellverteilung: $x_0 \sim \pi_0$.
    \item Setze die Schrittweite: $\Delta t = \frac{1}{N}$.
    \item Schleife für jeden Integrationsschritt $k = 0 \dots N-1$:
    \begin{align}
        t &= k \cdot \Delta t \\
        v &= v_\theta(x_t, t) \\
        x_{t + \Delta t} &= x_t + v \cdot \Delta t
    \end{align}
\end{enumerate}
\textbf{Output:} Das fertig generierte synthetische Datenobjekt $x_1$.

\section{Empirische Messergebnisse}
Die folgende Performancematrix wurde direkt aus der numerischen Live-Simulation extrahiert:

\begin{table}[h]
\centering
\caption{\textbf{Ergebnistabelle der Live-Messung}}
\vspace{0.2cm}
\begin{tabular}{lccc}
\toprule
\textbf{Algorithmus} & \textbf{NFE (Netzwerkaufrufe) $\downarrow$} & \textbf{Inferenz-Zeit (ms) $\downarrow$} & \textbf{Abweichung (MSE)} \\
\midrule
DDPM (Diffusion) & 10 & 12,18 & 1,8597 \\
Flow-Matching    & 10 & 14,29 & 2,2676 \\
\textbf{Rectified Flow} & \textbf{1} & \textbf{1,46} & \textbf{2,1518} \\
\bottomrule
\end{tabular}
\end{table}

\section{Wissenschaftliche Kommentierung der Forschungsfragen}

\subsection{Evaluierung von Ziel 1 (Wall-Clock Cost \& NFE-Reduktion)}
Die empirischen Daten bestätigen die theoretischen Kernvorteile von \textbf{Rectified Flow (RF)} hinsichtlich der Inferenz-Effizienz vollumfänglich. Wenn die Modelle an iterative Standard-Solver gebunden sind, benötigen DDPM und Flow-Matching eine identische Allokation von $\text{NFE} = 10$. Dies führt zu proportionalen Berechnungsverzögerungen auf der Hardware ($12,18\text{ ms}$ bzw. $14,29\text{ ms}$). 

Da Rectified Flow jedoch das Vektorfeld im Training mathematisch begradigt, entkoppelt es den Generierungsprozess vollständig von progressiven Integrationsschleifen. RF erreicht eine optimale \textbf{NFE = 1}, wodurch die reale Inferenz-Zeit (\textit{Wall-Clock Cost}) auf nur \textbf{1,46 ms} kollabiert. Dies entspricht einer \textbf{9,78-fachen Beschleunigung} gegenüber dem klassischen Flow-Matching auf identischer Architektur und beweist das Erreichen des ersten Forschungsziels für echtzeittaugliche Systeme.

\subsection{Evaluierung von Ziel 2 (Formulierungstabilität \& Linientreue)}
Die Stabilität der mathematischen Formulierung zeigt sich darin, wie resistent ein Framework gegen eine drastische Reduktion der Auflösungsschritte reagiert. Erzwingt man bei klassischen Diffusionsmodellen extrem niedrige Schrittketten, bricht die Generierungsqualität ein, da ihr stochastischer Rauschplan ($\beta_t$) stark diskretisiert wird und instabil fluktuiert. Flow-Matching stabilisiert dies, indem es den Rauschplan durch ein deterministisches ODE-Problem ersetzt, und erzielt in der Messung eine \textbf{Abweichung (MSE) von 2,2676}.

Betrachtet man nun \textbf{Rectified Flow bei einer NFE = 1}, erreicht das Modell einen \textbf{MSE von 2,1518}. Dass RF neun komplette Rechenschritte überspringen kann und dennoch einen \textit{geringeren} Abweichungsfehler aufweist als das iterative Flow-Matching, liefert den mathematischen Beweis für die überlegene Stabilität des unbeschränkten Least-Squares-Lernziels. Die Trajektorien werden im Training so präzise begradigt, dass ein einzelner linearer Schritt weniger Integrationsdrift erzeugt als das Navigieren durch eine gekrümmte, mehrstufige Kurve.

\end{document}